\documentclass[11pt,a4paper]{article}
\usepackage[margin=25mm,headheight=14pt]{geometry}
\usepackage{fontspec}
\setmainfont{lmroman10-regular.otf}[BoldFont=lmroman10-bold.otf,ItalicFont=lmroman10-italic.otf,BoldItalicFont=lmroman10-bolditalic.otf]
\setsansfont{lmsans10-regular.otf}
\setmonofont{lmmono10-regular.otf}[Scale=MatchLowercase]
\newfontfamily\authorsymbol{Noto Sans Symbols}[Scale=MatchUppercase]
\usepackage{amsmath,amssymb,amsthm,mathtools,booktabs}
\usepackage{microtype,needspace}
\usepackage[hidelinks,unicode]{hyperref}
\usepackage{fancyhdr}
\pagestyle{fancy}\fancyhf{}
\fancyhead[L]{\small\itshape Exact counting with bounded tap span}
\fancyhead[R]{\small\thepage}\renewcommand{\headrulewidth}{0pt}
\setlength{\parindent}{1.2em}\setlength{\parskip}{0pt}\setlength{\emergencystretch}{2em}
\newtheorem{theorem}{Theorem}
\newtheorem{lemma}[theorem]{Lemma}
\newtheorem{proposition}[theorem]{Proposition}
\newtheorem{corollary}[theorem]{Corollary}
\theoremstyle{definition}\newtheorem{definition}[theorem]{Definition}
\newcommand{\B}{\{0,1\}}
\newcommand{\maj}{\operatorname{Maj}}
\newcommand{\rot}{\operatorname{ROR}}
\newcommand{\shr}{\operatorname{SHR}}
\newcommand{\bits}[2]{#1_{#2}}
\title{\LARGE Exact Counting for a Boolean--Modular Equation\\with Bounded Tap Span}
\author{lmtlssss{\authorsymbol 🜐}}
\date{27 October 2026}
\hypersetup{pdftitle={Exact Counting for a Boolean–Modular Equation with Bounded Tap Span},pdfauthor={lmtlssss🜐},pdfsubject={Exact counting, finite-state algorithms and modular arithmetic},pdfkeywords={bounded frontier, modular arithmetic, exact counting, bit rotation, carry propagation}}
\begin{document}
\maketitle\thispagestyle{plain}
\begin{abstract}
Let $L_w(x)=\rot_{r_1}(x)\oplus\rot_{r_2}(x)\oplus\shr_s(x)$, where
$0<s<r_1<r_2<w$. We count all solutions of
$L_w(a+D)+L_w(a)=S$ modulo $2^w$ and recover one solution when the count is positive.
Set $h=r_2-s$. For each cyclic boundary, the last $h$ input bits and one entering
carry determine every legal continuation. The output carry follows from the target
and the input difference. An exact dynamic program therefore uses at most
$w2^{2h+2}$ nonterminal states. For fixed $h$, this count is linear in $w$;
straightforward arbitrary-precision arithmetic gives a quadratic bit-operation bound.
The proof includes cyclic closure and unique coverage. Exact implementations and
complete small-width comparisons accompany a 256-bit example with
$5{,}784{,}112{,}988{,}160$ solutions.
\end{abstract}
\noindent\textit{Keywords.} Exact counting; finite-state computation; modular arithmetic;
bit rotation; carry propagation.

\section{The equation and the parameter}
Identify a $w$-bit word with an integer in $\{0,\ldots,2^w-1\}$.
Bit $a_0$ is the least significant bit. Write $R_r=\rot_r$ for cyclic right
rotation and $H_s=\shr_s$ for logical right shift. Thus
\[
 (R_ra)_i=a_{(i+r)\bmod w},\qquad
 (H_sa)_i=\begin{cases}a_{i+s},&i+s<w,\\0,&i+s\geq w.\end{cases}
\]
Fix integers $0<s<r_1<r_2<w$ and define
\begin{equation}\label{eq:map}
 L_w=R_{r_1}\oplus R_{r_2}\oplus H_s,
 \qquad h=r_2-s,\qquad d=r_2-r_1.
\end{equation}
In particular, $1\leq d<h<w$. For supplied words $D,S$, the required set is
\begin{equation}\label{eq:fiber}
 \mathcal A(D,S)=\{a:L_w((a+D)\bmod 2^w)+L_w(a)\equiv S\pmod{2^w}\}.
\end{equation}
The output consists of $|\mathcal A(D,S)|$ and one member when the set is nonempty.
The linear map $L_w$ need not be invertible.

Finite-state methods for arithmetic with carries are established.
Mouha et al. developed the differential analysis of S-functions
\cite{mouha2011}; Velichkov et al. treated rotation within that framework
\cite{velichkov2011}. The contribution here is an explicit sufficient state
for \eqref{eq:fiber}, together with its boundary conditions and a parameterized
counting bound. The proof does not claim a new general dynamic-programming principle.

\Needspace{11\baselineskip}
\begin{theorem}\label{thm:main}
For the parameters in \eqref{eq:map}, there is an exact count-and-witness algorithm
for \eqref{eq:fiber} with at most
\begin{equation}\label{eq:bound}
 N(w,h)=w\,2^{2h+2}
\end{equation}
nonterminal state visits and at most $2N(w,h)$ candidate transitions.
Initialization requires $O(h2^{h+1}+w)$ elementary work. A direct implementation
uses $O(w^2 2^{2h+2})$ bit operations and $O(w^2 2^{h+1})$ bits of storage.
These bounds include exact counts and a witness, but not an explicit list of all solutions.
\end{theorem}

The exponent depends on the tap span $h$, not directly on $w$.
For fixed $h$, the number of states is linear in the word length.
If $h$ grows with $w$, the same formula retains its exponential dependence on $h$.
The proof of Theorem~\ref{thm:main} occupies Sections~2--5.

\section{The two carry equations}
Put $b=(a+D)\bmod 2^w$ and let $t_j$ be the carry into input position $j$.
For $0\leq j<w$,
\begin{align}
 t_0&=0, & b_j&=a_j\oplus D_j\oplus t_j,
 &t_{j+1}&=\maj(a_j,D_j,t_j).\label{eq:input}
\end{align}
Here $\maj(u,v,z)$ is one precisely when at least two arguments are one.
The final carry $t_w$ is discarded. Set $\delta=a\oplus b$. Then
\begin{equation}\label{eq:delta}
 \delta_j=D_j\oplus t_j.
\end{equation}
Write $x=L_w(a)$, $y=L_w(b)$, and let $c_i$ be the carry into position $i$
of the output addition. Its sum-bit equation is
\begin{equation}\label{eq:outputcarry}
 c_i=S_i\oplus x_i\oplus y_i
     =S_i\oplus L_w(\delta)_i.
\end{equation}
The equality on the right uses binary linearity of $L_w$.
The remaining output conditions are
\begin{equation}\label{eq:output}
 c_0=0,\qquad c_{i+1}=\maj(x_i,y_i,c_i)\quad(0\leq i<w-1).
\end{equation}
Again the carry out of position $w-1$ is discarded.

Process input positions in the order $j=0,\ldots,w-1$ and associate
\begin{equation}\label{eq:index}
 i=(j-r_2)\bmod w.
\end{equation}
For any word $v$, its output bit at this position is
\begin{equation}\label{eq:taps}
 L_w(v)_i=v_j\oplus v_{(j-d)\bmod w}
       \oplus \varepsilon_i v_{(j-h)\bmod w},\qquad
 \varepsilon_i=\mathbf 1_{\{i+s<w\}}.
\end{equation}
All noncurrent taps lie within the preceding $h$ cyclic input positions.
The factor $\varepsilon_i$ preserves the zero-filled right shift.

\begin{lemma}[Local transition]\label{lem:transition}
Suppose the preceding $h$ bits of $a,b$ and the entering input carry $t_j$ are known.
Then $c_i$ is known before $a_j$ is chosen. For either choice $a_j\in\B$,
the validity of the next transition is decidable from these data and $D,S$.
\end{lemma}
\begin{proof}
Equation~\eqref{eq:delta} gives $\delta_j$ without $a_j$.
Equations~\eqref{eq:outputcarry} and \eqref{eq:taps} then give $c_i$.
A choice of $a_j$ determines $b_j,t_{j+1},x_i,y_i$ by \eqref{eq:input}
and \eqref{eq:taps}. When $i<w-1$, compute the required next carry from
$S_{i+1}\oplus L_w(\delta)_{i+1}$ and compare it with $\maj(x_i,y_i,c_i)$.
The new history and $\delta_{j+1}=D_{j+1}\oplus t_{j+1}$ supply its taps.
At $j=w-1$, the next input index is zero; use $\delta_0=D_0$ because $t_0=0$.
At $i=w-1$, impose no outgoing output-carry condition. At $i=0$, impose $c_0=0$.
These rules give all sum-bit and carry conditions, including both discarded carries.
\end{proof}

\section{The cyclic boundary}
The first transition uses the high $h$ input bits.
Enumerate a boundary seed
\begin{equation}\label{eq:seed}
 \sigma=(\alpha,\tau)\in\B^h\times\B,
\end{equation}
where $\alpha=(a_{w-h},\ldots,a_{w-1})$ and $\tau=t_{w-h}$.
Add the known high $h$ bits of $D$ to $\alpha$, with entering carry $\tau$.
This determines a high-bit word $\beta=(b_{w-h},\ldots,b_{w-1})$.
Use $\alpha,\beta$ as the initial cyclic histories, while the actual forward
input carry starts at $t_0=0$.

At positions $j<w-h$, allow both values of $a_j$.
At positions $j\geq w-h$, require the corresponding seeded bit of $\alpha$.
Require $t_{w-h}=\tau$ at entry to the seeded tail, and require the produced
$b$-tail to equal $\beta$. Apply Lemma~\ref{lem:transition} throughout.

\begin{lemma}[Exact coverage]\label{lem:coverage}
Accepted paths, over all seeds \eqref{eq:seed}, are in bijection with
$\mathcal A(D,S)$.
\end{lemma}
\begin{proof}
An accepted path fixes one complete word $a$.
Its input additions and the tail-carry agreement give
$b=(a+D)\bmod 2^w$. The seeded histories agree with the completed tails.
Lemma~\ref{lem:transition} enforces every output sum bit and its carry condition.
Thus the word satisfies \eqref{eq:fiber}.

Conversely, a solution $a$ fixes its high $h$ bits and the unique carry $t_{w-h}$.
It therefore fixes exactly one seed. Its actual bits follow a legal path,
since each transition is an identity of the two additions.
Distinct words differ in a chosen bit; distinct seeds cannot describe the same word.
This proves both completeness and absence of duplicate counting.
\end{proof}

\section{A sufficient state of \texorpdfstring{$h+1$}{h+1} bits}
For a fixed seed and a fixed position $j\geq h$, define
\begin{equation}\label{eq:state}
 Q_j=\bigl(a_{j-h},\ldots,a_{j-1},\,t_{j-h}\bigr).
\end{equation}
This is a binary state of length $h+1$.

\begin{lemma}[Sufficient frontier]\label{lem:frontier}
For fixed $D,S$, seed $\sigma$, and position $j\geq h$, any two valid prefixes
with the same $Q_j$ have exactly the same legal suffixes.
\end{lemma}
\begin{proof}
The $h$ input bits in $Q_j$, the known $D$-segment and $t_{j-h}$ determine the
corresponding $h$ bits of $b$ and the carry $t_j$ by ordinary addition.
Equation~\eqref{eq:taps} uses only those histories and subsequent input bits.
Equation~\eqref{eq:outputcarry} determines the output carry from the same data.
The remaining seeded-tail restrictions are fixed by $\sigma$.
Consequently both prefixes admit the same next choices, with the same successor
states. Induction over the remaining positions gives identical legal suffixes.
\end{proof}

For $j<h$, the initial tail is fixed, and the $h$-bit input history contains
all $j$ selected low bits. It identifies the prefix uniquely.
There are at most $2^j\leq2^{h+1}$ states at these positions as well.
If an implementation retains both input histories, it can recover the entering
history carry directly:
\begin{equation}\label{eq:oldcarry}
 t_{j-h}=a_{j-h}\oplus b_{j-h}\oplus D_{j-h}.
\end{equation}
The additional history is therefore redundant, rather than an additional
independent component of the state.

\section{Counting, witnesses and cost}
For each seed, form the layered directed graph of legal transitions.
Merge equal states within a layer according to Lemma~\ref{lem:frontier}.
Attach weight one to the initial state. For each valid transition, add the
source weight to its destination weight. The sum of terminal weights is the
number of accepted words for that seed. Add the totals over all seeds.

A representative prefix can accompany each positive weight.
When two prefixes merge, retain either representative: the frontier lemma gives
both the same suffix set. A positive terminal weight therefore supplies a
witness. A suffix-memoized implementation is equivalent. It stores the number
of legal suffixes and one suffix witness at each state.

\begin{proof}[Proof of Theorem~\ref{thm:main}]
Lemma~\ref{lem:coverage} identifies accepted paths with the required solutions.
The frontier lemma justifies every merge. Weighted path counting and witness
reconstruction therefore give the stated output exactly.
There are $2^{h+1}$ seeds, $w$ nonterminal layers, and at most $2^{h+1}$ states
per seed and layer. Their product is \eqref{eq:bound}.
Each state has at most two candidate transitions.
Preparing the high-bit addition for every seed costs $O(h2^{h+1})$;
reading the input bits costs $O(w)$.

Counts never exceed $2^w$, so $w+1$ bits suffice.
A representative word needs $w$ bits. The local transition uses a bounded
number of bit operations if the redundant histories are retained.
Conservative indexing and arbitrary-precision addition cost $O(w)$ per
state or transition. This gives $O(w^2 2^{2h+2})$ bit operations.
A suffix table for one seed has at most $w2^{h+1}$ entries of $O(w)$ bits;
reuse the table for the next seed. This gives the stated storage bound.
\end{proof}

\begin{corollary}
For every fixed tap span $h$, the count-and-witness problem \eqref{eq:fiber}
has polynomial bit complexity in $w$, independently of the size of its fiber.
\end{corollary}
Explicit enumeration has a different output contract: writing $N$ distinct
$w$-bit words requires $\Omega(wN)$ output bits.
Likewise, an added condition can distinguish prefixes that \eqref{eq:state}
merges. Such a condition requires a sufficient state of its own.
The theorem concerns \eqref{eq:fiber} with the stated taps and boundary rules.

\section{A complete 256-bit example}
Use $(r_1,r_2,s)=(17,19,10)$, so $h=9$.
For reproducibility, the following hexadecimal strings are the supplied words
$D$ and $S$. Concatenate each pair of lines, with the upper line most significant.
\begin{center}\small\ttfamily
\begin{tabular}{cl}
$D$ & bd8cc0534cb95939a2e5cbed5a8e57cc\\[-1pt]
    & 809ec58eeee0e0317a2d380ff7a92156\\[3pt]
$S$ & 9cec8b4181dcf3012fa4710a9b9a8ece\\[-1pt]
    & dc001740cc795ef2b57d4bad916262f9
\end{tabular}
\end{center}
The exact result is
\begin{equation}\label{eq:example}
 |\mathcal A(D,S)|=5{,}784{,}112{,}988{,}160.
\end{equation}
The native suffix implementation visits $558{,}622$ nonterminal states and
examines $1{,}115{,}268$ transitions. Its $1{,}024$ boundary seeds cover every
possible high nine-bit word and entering carry. These are deterministic
execution counts for this instance, not a bound for all inputs.
The general bound here is $256\cdot2^{20}=268{,}435{,}456$ states.
The full candidate domain has $2^{256}$ words.

Two Python implementations reproduce \eqref{eq:example}. One uses the canonical
frontier in forward order. The other retains both input histories and an
explicit output carry, including its cyclic boundary. Their agreement checks
the eliminated-carry construction through a different state representation.
Each returned witness also satisfies \eqref{eq:fiber} by direct word arithmetic.

\section{Verification and available sources}
The accompanying code contains the integer C++ suffix implementation, a canonical
Python forward implementation, and a Python reference with explicit output carries.
It also contains the fixed inputs, exact results and the figure construction data.
The verification covers the following complete paths.
\begin{center}
\begin{tabular}{lr}
\toprule
Comparison & Queries\\
\midrule
All $D,S$ at widths $4$--$7$, declared tap sets & 21,760\\
Every tap triple at widths $4$--$8$, fixed test inputs & 840\\
Three implementations on fixed widths $24$--$256$ & 54\\
\bottomrule
\end{tabular}
\end{center}
The first comparison builds the complete forward histogram for each $D$ and
checks each target count. The second includes zero, all-one and independently
selected inputs. The third includes empty and nonempty fibers, with counts
represented as arbitrary-precision integers. All comparisons agree.
These executions supplement the proof; the general theorem follows from
Lemmas~\ref{lem:transition}--\ref{lem:frontier} and the state count.

The implementations enforce declared resource limits. An interrupted search
returns an incomplete result, rather than an empty fiber.
The state bound specifies combinatorial work; the bit-operation bound also charges exact integer arithmetic.
The specific sufficient state and carry elimination are the mathematical content;
finite-state arithmetic counting itself has the prior literature cited above.

\noindent\textit{Reproducibility.}
The versioned release supplies the full code, exact example, verification receipt,
editable figure and manuscript source. Its DOI is \texttt{10.5281/zenodo.23697063}.

\begin{thebibliography}{9}
\bibitem{mouha2011}
N. Mouha, V. Velichkov, C. De Canni\`ere and B. Preneel.
\emph{The Differential Analysis of S-Functions}.
In Selected Areas in Cryptography, LNCS 6544, pp. 36--56, Springer, 2011.
\href{https://doi.org/10.1007/978-3-642-19574-7_3}{doi:10.1007/978-3-642-19574-7\_3}.
\bibitem{velichkov2011}
V. Velichkov, N. Mouha, C. De Canni\`ere and B. Preneel.
\emph{The Additive Differential Probability of ARX}.
In Fast Software Encryption, LNCS 6733, pp. 342--358, Springer, 2011.
\href{https://doi.org/10.1007/978-3-642-21702-9_20}{doi:10.1007/978-3-642-21702-9\_20}.
\end{thebibliography}
\end{document}
